\begin{lemma}[Surgery induction]
  \label{surgery-induction}
  Let $E$ be a metric space equipped with its Borel $\sigma$-algebra, let $GP$ be a geodesic
  pairing on $E$ (\noderef{GeodesicPairing}), let $\gamma \in \Theta(E)$ (\noderef{Curve}) be a
  closed (\noderef{IsClosedCurve}), piecewise-geodesic (\noderef{IsPiecewiseGeodesic}) curve, let
  $\epsilon,\delta \in \mathbb{R}$ with $0<\epsilon<1$ and $0<\delta$, and let $n \in \mathbb{N}$
  with $n>0$. Then there exist $N \in \mathbb{N}$ with $N \ge 1$ and closed, piecewise-geodesic
  curves $(g_j)_{j=1}^N$ such that
  \begin{align}
    \label{eq:SurgeryIndCurrent}
    [[\gamma]](\omega) &= \sum_{j=1}^N [[g_j]](\omega) \qquad\text{for every } \omega \in
    \mathcal{D}^1(E) \quad\text{(\noderef{curveCurrent})}, \\
    \label{eq:SurgeryIndLength}
    \sum_{j=1}^N \length(g_j) &\le \mathrm{surgeryPotential}(\epsilon,\delta,n,\length(\gamma),
    \nsmallpieces(\gamma,\delta)) \quad\text{(\noderef{surgeryPotential}, \noderef{smallPieceCount})},\\
    \label{eq:SurgeryIndMorrey}
    \Morrey{\mu_{g_j}} &\le 4\epsilon^{-1}+2n+10 \qquad\text{for every } j
    \quad\text{(\noderef{curveMeasure}, \noderef{morreyNorm})}.
  \end{align}

  This is the content of paper Lemma~3.1 (the paper's own label \texttt{surgery-lem}, paper.tex
  line 597), specialized to $L:=\length(\gamma)$, $m:=\nsmallpieces(\gamma,\delta)$
  \emph{before} the algorithm of paper.tex lines 1005-1075 (``Algorithm 1'') has been run: paper
  Lemma~3.1 itself, with the paper's literal bound $\bigl(1+\tfrac{2\epsilon}{1-\epsilon}+
  \tfrac{12\epsilon^{-1}}{n(1-\epsilon)}\bigr)\length(\gamma)$, will be recovered from this node
  by a planned follow-up node, using \noderef{InitialDeltaExists} to arrange $m=0$ at the start
  (that node is outside the present scope; this node supplies its entire mathematical content).
  Below, the whole of Algorithm~1 -- its counters $i,j,T_1,T_2$ and its
  stored sequence $(\beta_j)$ -- is replaced by a single strong induction on a natural-number
  termination measure; the case split of the loop body becomes the case split on
  $\Xden(\delta,\epsilon,n)$ versus $\Xlsi(\delta,\epsilon)$ membership below, and the algorithm's
  own two-stage bound $T_I \le \length(\gamma)/((1-\epsilon)\delta)$, $T_{II}\le(3/n)T_I$ (paper.tex
  1058-1062) is replaced by the single potential inequality \eqref{eq:SurgeryIndLength}, established
  directly by an amortized argument (the length steps proved below) rather
  than by counting cuts of each type separately.

  \paragraph{Why $n>0$ is mandatory, and is a faithful hypothesis, not a deviation.} The
  coefficient $12\epsilon^{-1}/(n(1-\epsilon))$ appearing in \noderef{surgeryPotential} is literally
  undefined at $n=0$ (division by zero in a context where the paper intends a genuine positive
  quantity), so \eqref{eq:SurgeryIndLength} cannot even be stated at $n=0$ without a convention. More
  substantively, the induction's Type~I termination step (proved below) uses
  $n>0$ essentially: the measure drops by exactly $n$ at each Type~I cut, and at $n=0$ this drop
  vanishes, so the induction need not terminate. This matches the paper's own use of $n$: the
  standing hypothesis of the entire surgery construction (Def.~3.5, paper.tex, restated at the head
  of paper Lemma~3.1) is $n \in \mathbb{N}$ with no explicit positivity clause stated, but every
  citing use of paper Lemma~3.1 (paper Cor.~3.2, \texttt{surgery-eta}, paper.tex line 628, with
  $n = \lceil \epsilon^{-2}\rceil \ge 1$) supplies $n \ge 1$, so requiring $n>0$ here changes nothing about
  which instances of the lemma are actually used.
\end{lemma}
\begin{proof}
  \paragraph{The termination measure.} For a closed, piecewise-geodesic curve $\eta$, define the
  natural number
  \[
    M(\eta) := (n+3) \cdot \left\lfloor \frac{\length(\eta)}{(1-\epsilon)\delta} \right\rfloor +
    \nsmallpieces(\eta,\delta)
    \quad \in \mathbb{N}
    ,
  \]
  where $\lfloor\cdot\rfloor$ is the floor of a nonnegative real number (well defined since
  $\length(\eta)\ge0$ by \noderef{Curve} and $(1-\epsilon)\delta>0$, as $0<\epsilon<1$ and
  $0<\delta$). Write $\Phi(L,m) := \mathrm{surgeryPotential}(\epsilon,\delta,n,L,m)$
  (\noderef{surgeryPotential}) for the potential of the statement, with its defining formula
  \[
    \Phi(L,m) = A \cdot L + C \cdot m
    , \qquad
    A := 1 + \frac{2\epsilon}{1-\epsilon} + \frac{12\epsilon^{-1}}{n(1-\epsilon)}
    , \qquad
    C := \frac{4\epsilon^{-1}\delta}{n}
    .
  \]
  Since $0<\epsilon<1$ and $n>0$, all four summands $1,2\epsilon/(1-\epsilon),
  12\epsilon^{-1}/(n(1-\epsilon))$ are nonnegative, so
  \begin{equation}
    \label{eq:A-ge-one}
    A \ge 1
    , \qquad\text{hence}\qquad A - 1 \ge 0
    ,
  \end{equation}
  and, since additionally $\delta>0$, $C>0$. A direct computation (clearing denominators, using
  $1-\epsilon\ne0$ and $n\ne0$) gives the two algebraic identities used below:
  \begin{equation}
    \label{eq:A-minus-one-identity}
    (A-1)(1-\epsilon) = 2\epsilon + \frac{12\epsilon^{-1}}{n}
    , \qquad\qquad
    C \cdot n = 4\epsilon^{-1}\delta
    .
  \end{equation}

  \paragraph{Statement of the induction.} We prove, for every closed piecewise-geodesic curve
  $\eta$, the claim $Q(\eta)$: ``there exist $N\ge1$ and closed piecewise-geodesic $(g_j)_{j=1}^N$
  satisfying \eqref{eq:SurgeryIndCurrent}--\eqref{eq:SurgeryIndMorrey} with $\gamma$ replaced by
  $\eta$ throughout (and $\length(\gamma),\nsmallpieces(\gamma,\delta)$ replaced by
  $\length(\eta),\nsmallpieces(\eta,\delta)$ in \eqref{eq:SurgeryIndLength})''. We prove $Q(\eta)$
  for every such $\eta$ by strong induction on $M(\eta) \in \mathbb{N}$: fix $\eta$ closed and
  piecewise-geodesic, and assume the induction hypothesis (IH) that $Q(\eta')$ holds for every
  closed piecewise-geodesic $\eta'$ with $M(\eta') < M(\eta)$. Applying the resulting $Q(\gamma)$
  to the $\gamma$ of the lemma's statement (closed and piecewise-geodesic by hypothesis) then gives
  exactly the lemma's conclusion.

  By the law of excluded middle applied twice, exactly one of the following three cases holds for
  $\eta$: (A) $\eta \notin \Xden(\delta,\epsilon,n)$; (B) $\eta \in \Xden(\delta,\epsilon,n)$ and
  $\eta \in \Xlsi(\delta,\epsilon)$; (C) $\eta \in \Xden(\delta,\epsilon,n)$ and $\eta \notin
  \Xlsi(\delta,\epsilon)$ (\noderef{IsDeltaEpsN}, \noderef{IsLSI}).

  \paragraph{A finite-sum splitting identity, used in both cut cases.} If $N'\in\mathbb{N}$, $d$ is
  a curve, and $(h_i)_{i=1}^{N'}$ is a family of curves, then setting $N:=N'+1$ and $g_1:=d$,
  $g_{i+1}:=h_i$ for $1\le i\le N'$ (i.e.\ $(g_j)_{j=1}^N$ is $d$ prepended to $(h_i)_{i=1}^{N'}$),
  for any function $F$ of a curve valued in an additive commutative monoid,
  \begin{equation}
    \label{eq:sum-split}
    \sum_{j=1}^N F(g_j) = F(d) + \sum_{i=1}^{N'} F(h_i)
    ,
  \end{equation}
  simply by separating the $j=1$ term (which is $F(g_1)=F(d)$) from the remaining $N'$ terms
  (which are $F(g_{i+1})=F(h_i)$ for $i=1,\dots,N'$, a relabelling of the sum's index). This is
  used below with $F$ ranging over $\omega\mapsto[[\,\cdot\,]](\omega)$ (additive in
  $\mathbb{R}$) and $\cdot\mapsto\length(\cdot)$ (additive in $\mathbb{R}$), and separately to
  transport the pointwise bound $F\le K$ over $(g_j)_j$ from $F(d)\le K$ and $F(h_i)\le K$ for all
  $i$ (immediate, since $j=1$ gives $d$ and $j=i+1$ gives $h_i$).

  \paragraph{Case (A): $\eta \notin \Xden(\delta,\epsilon,n)$.} By \noderef{CutTypeII} applied to
  $GP,\eta,\delta,\epsilon,n$ (valid: $0<\epsilon<1$, $0<\delta$, $\eta$ closed and
  piecewise-geodesic, and this case's hypothesis), there exist $\eta',d \in \Theta(E)$, both closed
  and piecewise-geodesic, such that
  \[
    [[\eta]](\omega)=[[\eta']](\omega)+[[d]](\omega) \ (\forall\omega),\quad
    \Morrey{\mu_d} \le 2(n+2\epsilon^{-1}+2),\quad
    \length(\eta')\le\length(\eta),
  \]
  \[
    \length(\eta')+\length(d) \le \length(\eta)+4\epsilon^{-1}\delta
    , \qquad
    \nsmallpieces(\eta',\delta)+n \le \nsmallpieces(\eta,\delta)
    .
  \]

  \emph{Measure drop.} Write $m:=\nsmallpieces(\eta,\delta)$, $m':=\nsmallpieces(\eta',\delta)$,
  $F:=\lfloor\length(\eta)/((1-\epsilon)\delta)\rfloor$, $F':=\lfloor\length(\eta')/((1-\epsilon)
  \delta)\rfloor$. Since $\length(\eta')\le\length(\eta)$ and dividing by the positive number
  $(1-\epsilon)\delta$ preserves order, $\length(\eta')/((1-\epsilon)\delta) \le
  \length(\eta)/((1-\epsilon)\delta)$; since the floor function is monotone on the reals,
  $F'\le F$, hence $(n+3)F' \le (n+3)F$. Since $m'+n\le m$ and $n>0$, $m' < m$ (indeed $m'\le m-n$).
  Combining, $M(\eta')=(n+3)F'+m' \le (n+3)F+m' \le (n+3)F+(m-n) = M(\eta)-n < M(\eta)$, using
  $n>0$ for the last strict inequality. So
  \begin{equation}
    \label{eq:typeII-measure-drop}
    M(\eta') < M(\eta)
    .
  \end{equation}

  \emph{The length step.} We claim
  \begin{equation}
    \label{eq:typeII-arith}
    \length(d) + \Phi(\length(\eta'),m') \le \Phi(\length(\eta),m)
    .
  \end{equation}
  Indeed, $\Phi(\length(\eta),m)-\Phi(\length(\eta'),m') = A\bigl(\length(\eta)-\length(\eta')\bigr)
  + C(m-m')$. Since $\length(\eta)-\length(\eta')\ge0$ (shown above) and $A-1\ge0$
  \eqref{eq:A-ge-one}, $(A-1)(\length(\eta)-\length(\eta'))\ge0$, so $A(\length(\eta)-\length(\eta'))
  \ge \length(\eta)-\length(\eta')$. Since $m-m'\ge n$ and $C>0$, $C(m-m')\ge Cn=4\epsilon^{-1}
  \delta$ by \eqref{eq:A-minus-one-identity}. Hence
  \[
    \Phi(\length(\eta),m)-\Phi(\length(\eta'),m') \ge \bigl(\length(\eta)-\length(\eta')\bigr) +
    4\epsilon^{-1}\delta \ge \bigl(\length(\eta)-\length(\eta')\bigr) +
    \bigl(\length(\eta')+\length(d)-\length(\eta)\bigr) = \length(d)
    ,
  \]
  the middle inequality from $\length(\eta')+\length(d)\le\length(\eta)+4\epsilon^{-1}\delta$
  rearranged. This is \eqref{eq:typeII-arith}.

  \emph{Assembly.} By \eqref{eq:typeII-measure-drop} the IH applies to $\eta'$ (closed and
  piecewise-geodesic, as produced by \noderef{CutTypeII}): there exist $N'\ge1$ and closed
  piecewise-geodesic $(h_i)_{i=1}^{N'}$ with $[[\eta']](\omega)=\sum_i[[h_i]](\omega)$,
  $\sum_i\length(h_i)\le\Phi(\length(\eta'),m')$, and $\Morrey{\mu_{h_i}}\le4\epsilon^{-1}+2n+10$
  for all $i$. Set $N:=N'+1$ and $(g_j)_{j=1}^N$ the family obtained by prepending $d$ to
  $(h_i)_{i=1}^{N'}$, as in the splitting identity above. Then $N\ge1$; each $g_j$ is closed and
  piecewise-geodesic ($g_1=d$ from \noderef{CutTypeII}, $g_{i+1}=h_i$ from the IH). By
  \eqref{eq:sum-split} with $F(\cdot)=[[\cdot]](\omega)$, $[[\eta]](\omega) = [[\eta']](\omega)+
  [[d]](\omega) = \bigl(\sum_i[[h_i]](\omega)\bigr)+[[d]](\omega) = \sum_j[[g_j]](\omega)$, which is
  \eqref{eq:SurgeryIndCurrent} for $\eta$. By \eqref{eq:sum-split} with $F=\length$,
  $\sum_j\length(g_j) = \length(d)+\sum_i\length(h_i) \le \length(d)+\Phi(\length(\eta'),m') \le
  \Phi(\length(\eta),m)$, the last step by \eqref{eq:typeII-arith}; this is
  \eqref{eq:SurgeryIndLength} for $\eta$. Finally, for the Morrey bound: $g_1=d$ satisfies
  $\Morrey{\mu_d}\le2(n+2\epsilon^{-1}+2) = 4\epsilon^{-1}+2n+4 \le 4\epsilon^{-1}+2n+10$ (as reals,
  $4\le10$), and applying the (extended-real) monotonicity of the Morrey norm bound under this
  numeral inequality gives $\Morrey{\mu_d}\le4\epsilon^{-1}+2n+10$; each $g_{i+1}=h_i$ satisfies
  $\Morrey{\mu_{h_i}}\le4\epsilon^{-1}+2n+10$ directly from the IH. This establishes
  \eqref{eq:SurgeryIndMorrey} for $\eta$ and completes Case~(A).

  \paragraph{Case (B): $\eta\in\Xden(\delta,\epsilon,n)$ and $\eta\in\Xlsi(\delta,\epsilon)$.} Take
  $N:=1$ and $g_1:=\eta$ (the terminal step: the algorithm halts and outputs the current curve
  unchanged, paper.tex line 1007, ``\textsc{if} $\gamma_i$ is $(\delta,\epsilon)$-l.s.i.\
  \textsc{then} \ldots \textsc{return}''). Then $N\ge1$ trivially, and $g_1=\eta$ is closed and
  piecewise-geodesic by $\eta$'s own standing properties. The sum in
  \eqref{eq:SurgeryIndCurrent} has a single term, giving $[[\eta]](\omega)=[[g_1]](\omega)$
  trivially.

  \emph{The length bound.} We claim $\length(\eta) \le
  \Phi(\length(\eta),m)$ where $m:=\nsmallpieces(\eta,\delta)$. Indeed $\Phi(\length(\eta),m) = A
  \cdot\length(\eta)+C\cdot m \ge 1\cdot\length(\eta)+0 = \length(\eta)$, using $A\ge1$
  \eqref{eq:A-ge-one} together with $\length(\eta)\ge0$ (\noderef{Curve}) to get $A\cdot\length(\eta)
  \ge\length(\eta)$, and $C>0$, $m\ge0$ to get $C\cdot m\ge0$. This is exactly
  \eqref{eq:SurgeryIndLength}'s single-term instance, $\length(g_1)=\length(\eta) \le
  \Phi(\length(\eta),m)$.

  \emph{The Morrey bound.} Since $\eta\in\Xden(\delta,\epsilon,n)$ and $\eta\in\Xlsi(\delta,
  \epsilon)$, \noderef{FullBallLem} applies to $\eta,\delta,\epsilon,n$ (valid: $0<\delta$,
  $0<\epsilon<1$, and these two membership facts) and gives $\Morrey{\mu_\eta} \le
  4\epsilon^{-1}+2n+4 \le 4\epsilon^{-1}+2n+10$ (as reals, $4\le10$; converting via the monotonicity
  of the extended-real bound as in Case~(A)). This is \eqref{eq:SurgeryIndMorrey}'s single-term
  instance and completes Case~(B).

  \paragraph{Case (C): $\eta\in\Xden(\delta,\epsilon,n)$ and $\eta\notin\Xlsi(\delta,\epsilon)$.}
  By \noderef{CutTypeI} applied to $GP,\eta,\delta,\epsilon,n$ (valid: $0<\epsilon<1$, $0<\delta$,
  $\eta$ closed and piecewise-geodesic, and this case's two membership facts), there exist
  $\eta',d\in\Theta(E)$, both closed and piecewise-geodesic, and $\beta\in\mathbb{R}$, such that
  \[
    \delta\le\beta\le\tfrac12\length(\eta),\qquad
    [[\eta]](\omega)=[[\eta']](\omega)+[[d]](\omega)\ (\forall\omega),\qquad
    \Morrey{\mu_d}\le4\epsilon^{-1}+2n+10,
  \]
  \[
    \length(\eta')\le\length(\eta)-(1-\epsilon)\beta,\qquad
    \length(\eta')+\length(d)\le\length(\eta)+2\epsilon\beta,\qquad
    \nsmallpieces(\eta',\delta)\le\nsmallpieces(\eta,\delta)+3
    .
  \]
  (Unlike Case~(A), the Morrey bound on $d$ is \emph{already} exactly the target constant
  $4\epsilon^{-1}+2n+10$; no numeral adjustment is needed for this case.)

  \emph{Measure drop.} Write $m:=\nsmallpieces(\eta,\delta)$,
  $m':=\nsmallpieces(\eta',\delta)$, $F,F'$ as in Case~(A). Since $\delta\le\beta$ and $1-\epsilon>0$,
  $(1-\epsilon)\delta\le(1-\epsilon)\beta$, so $\length(\eta')\le\length(\eta)-(1-\epsilon)\beta\le
  \length(\eta)-(1-\epsilon)\delta$, i.e.\ $\length(\eta)-\length(\eta')\ge(1-\epsilon)\delta>0$.
  Writing $c:=(1-\epsilon)\delta>0$: since $\length(\eta')/c \le \length(\eta)/c - 1$ (dividing the
  length gap $\ge c$ by $c$), and $\lfloor\length(\eta')/c\rfloor\le\length(\eta')/c$, we get
  $\lfloor\length(\eta')/c\rfloor+1 \le \length(\eta)/c$, and since the left side is an integer,
  $\lfloor\length(\eta')/c\rfloor+1 \le \lfloor\length(\eta)/c\rfloor$, i.e.\ $F'+1\le F$. Hence
  $(n+3)F'+(n+3) \le (n+3)F$. Together with $m'\le m+3$ (i.e.\ $m' - 3\le m$, reorganized as
  $m'\le m+3$ used directly below):
  \[
    M(\eta') = (n+3)F'+m' \le \bigl((n+3)F-(n+3)\bigr) + (m+3) = M(\eta) - (n+3)+3 = M(\eta) - n <
    M(\eta)
    ,
  \]
  using $n>0$ for the final strict inequality. So $M(\eta')<M(\eta)$.

  \emph{The length step.} We claim
  \begin{equation}
    \label{eq:typeI-arith}
    \length(d) + \Phi(\length(\eta'),m') \le \Phi(\length(\eta),m)
    .
  \end{equation}
  From $\length(\eta')\le\length(\eta)-(1-\epsilon)\beta$, $\length(\eta)-\length(\eta') \ge
  (1-\epsilon)\beta$; since $A-1\ge0$ \eqref{eq:A-ge-one}, $(A-1)(\length(\eta)-\length(\eta')) \ge
  (A-1)(1-\epsilon)\beta = \bigl(2\epsilon+12\epsilon^{-1}/n\bigr)\beta$ by
  \eqref{eq:A-minus-one-identity}. Also, from $m'\le m+3$, $m-m'\ge-3$, so (using $C>0$)
  $C(m-m') \ge -3C = -12\epsilon^{-1}\delta/n$ (using $C=4\epsilon^{-1}\delta/n$). Combining,
  \[
    \Phi(\length(\eta),m)-\Phi(\length(\eta'),m') = \bigl(\length(\eta)-\length(\eta')\bigr) +
    (A-1)\bigl(\length(\eta)-\length(\eta')\bigr) + C(m-m')
  \]
  \[
    \ge \bigl(\length(\eta)-\length(\eta')\bigr) + 2\epsilon\beta + \frac{12\epsilon^{-1}}{n}\beta -
    \frac{12\epsilon^{-1}}{n}\delta
    = \bigl(\length(\eta)-\length(\eta')\bigr) + 2\epsilon\beta +
    \frac{12\epsilon^{-1}}{n}(\beta-\delta)
    .
  \]
  Since $\beta\ge\delta$ (this is the \emph{second}, independent use of $\beta\ge\delta$ in this
  proof: the first controlled the measure drop above, and this one absorbs the $+3$ small-edge
  slack of \eqref{eq:typeI-arith}'s hypothesis $m'\le m+3$ into a nonnegative correction term) and
  $12\epsilon^{-1}/n>0$, the last summand is $\ge0$, so
  \[
    \Phi(\length(\eta),m)-\Phi(\length(\eta'),m') \ge \bigl(\length(\eta)-\length(\eta')\bigr) +
    2\epsilon\beta \ge \length(d)
    ,
  \]
  the last inequality from $\length(\eta')+\length(d) \le \length(\eta)+2\epsilon\beta$ rearranged
  to $\length(d) \le (\length(\eta)-\length(\eta'))+2\epsilon\beta$. This is
  \eqref{eq:typeI-arith}.

  \emph{Assembly.} Exactly as in Case~(A): by the measure drop, the IH applies to $\eta'$, giving
  $N'\ge1$ and closed piecewise-geodesic $(h_i)_{i=1}^{N'}$ satisfying
  \eqref{eq:SurgeryIndCurrent}--\eqref{eq:SurgeryIndMorrey} for $\eta'$. Set $N:=N'+1$ and
  $(g_j)_{j=1}^N$ the family obtained by prepending $d$ to $(h_i)_{i=1}^{N'}$. Then $N\ge1$; each
  $g_j$ is closed and piecewise-geodesic ($g_1=d$ from \noderef{CutTypeI}, $g_{i+1}=h_i$ from the
  IH). The current identity and length bound follow exactly as in Case~(A), using
  \eqref{eq:sum-split} and \eqref{eq:typeI-arith} in place of \eqref{eq:typeII-arith}. For the
  Morrey bound, $g_1=d$ satisfies $\Morrey{\mu_d}\le4\epsilon^{-1}+2n+10$ directly (no numeral
  adjustment needed here, unlike Case~(A)), and each $g_{i+1}=h_i$ satisfies the bound from the IH.
  This completes Case~(C).

  \paragraph{Conclusion.} Cases (A), (B), (C) are exhaustive and each establishes $Q(\eta)$ from the
  IH (with, in Cases (A) and (C), the IH invoked at a curve $\eta'$ with strictly smaller measure
  $M(\eta')<M(\eta)$, and in Case (B) no recursive use of the IH at all). By strong induction on
  $M(\eta)\in\mathbb{N}$, $Q(\eta)$ holds for every closed piecewise-geodesic curve $\eta$. Applying
  this to $\eta:=\gamma$ (closed and piecewise-geodesic by hypothesis) gives the lemma's conclusion.
\end{proof}
